\begin{pastexam}[label={px:df3-109-ds-02}]{地方三等資訊處理 109 年 第二題}
（申論；20 分）

\begin{itemize}[leftmargin=2.2em]
\item[㈠] 請描述 order 為 m 的 B-tree 之特性。（6分）
\item[㈡] 請問 order 為 m 高度為 h 的 B-tree：(1)最多有幾個節點？最多有幾個 Key？（6分）(2)最少有幾個節點？
最少有幾個 Key？（8分）
\end{itemize}
\end{pastexam}

\begin{pastsoln}
答案：令 $c = \lceil m/2 \rceil$，高度 $h$ 以層數計（只有根時 $h = 1$）。 (1) 最多 $\dfrac{m^h - 1}{m - 1}$ 個節點、
$m^h - 1$ 個 key；(2) 最少 $1 + \dfrac{2(c^{h-1} - 1)}{c - 1}$ 個節點、$2c^{h-1} - 1$ 個 key。

\textbf{㈠} order 為 $m$ 的 B-tree（資結書定義~11.2，$m \ge 3$）是一棵 $m$ 路搜尋樹，滿足：
\begin{enumerate}[label=(\arabic*)]
\item 每個節點最多 $m$ 個子節點，也就是最多 $m - 1$ 個 key；節點內的 key 由小到大排列。
\item 根若不是葉節點，至少有 $2$ 個子節點（根至少 $1$ 個 key）。
\item 根以外的內部節點至少有 $\lceil m/2 \rceil$ 個子節點；根以外的每個節點至少有 $\lceil m/2 \rceil - 1$ 個 key。
\item 有 $k$ 個 key 的內部節點恰有 $k + 1$ 個子節點，第 $i$ 棵子樹的 key 都介於第 $i - 1$ 和第 $i$ 個 key 之間
（搜尋樹的順序）。
\item 所有葉節點都在同一層（樹永遠是平衡的）。
\end{enumerate}
因此 $n$ 個 key 的 B-tree 高度是 $\Theta(\log_m n)$，搜尋、插入、刪除只要讀 $O(\log_m n)$ 個節點，適合放在磁碟上
做索引。 插入時節點溢位就把中間的 key 往上送、分裂成兩個節點（見歷屆試題~\ref{px:gk3-114-ds-01}）；刪除時
key 太少就向兄弟借或和兄弟合併。

\textbf{㈡ (1) 最多}。 每個節點都放滿：$m$ 個子節點、$m - 1$ 個 key。 第 $i$ 層最多 $m^{i-1}$ 個節點，
\[
\text{節點數} \le 1 + m + m^2 + \cdots + m^{h-1} = \frac{m^h - 1}{m - 1},
\qquad
\text{key 數} \le (m - 1) \cdot \frac{m^h - 1}{m - 1} = m^h - 1.
\]

\textbf{(2) 最少}。 每個節點都放得最少：根 $1$ 個 key、$2$ 個子節點；其他內部節點 $c$ 個子節點；根以外的節點
$c - 1$ 個 key。 第 $1$ 層 $1$ 個節點，第 $i$ 層（$2 \le i \le h$）$2c^{i-2}$ 個節點：
\[
\text{節點數} \ge 1 + 2(1 + c + \cdots + c^{h-2}) = 1 + \frac{2(c^{h-1} - 1)}{c - 1},
\]
\[
\text{key 數} \ge 1 + (c - 1) \cdot \frac{2(c^{h-1} - 1)}{c - 1} = 2c^{h-1} - 1.
\]
這兩個下界都達得到（每個節點都取最少的那棵樹），是資結書定理~11.1 的下界。 $h = 1$ 時兩式都是 $1$（只有根，
至少 $1$ 個 key）。

\textbf{例子}。 $m = 5$（$c = 3$）、$h = 3$：最多 $1 + 5 + 25 = 31$ 個節點、$5^3 - 1 = 124$ 個 key；最少
$1 + 2 + 6 = 9$ 個節點、$2 \cdot 3^2 - 1 = 17$ 個 key。 $m = 3$（2-3 樹，$c = 2$）時最少的情形是每個節點只有
$1$ 個 key 的滿二元樹，$2^h - 1$ 個節點、$2^h - 1$ 個 key，和公式一致。

若題目的高度是數邊（只有根時 $h = 0$），把上面的 $h$ 換成 $h + 1$：最多 $\frac{m^{h+1} - 1}{m - 1}$ 個節點、
$m^{h+1} - 1$ 個 key；最少 $1 + \frac{2(c^{h} - 1)}{c - 1}$ 個節點、$2c^{h} - 1$ 個 key。
\end{pastsoln}
